\section{Следующие кольца нётеровы: а) кольцо главных идеалов; б) $Z[
\sqrt{d}]$ (d — любое целое число);
в) $Q[x1, . . . , x_n]$.}

\Def Область целостности $K$ называется \textit{нетёровым кольцом}, если выполнено одно из следующих условий:
\begin{itemize}
    \item [1] Если $I_0 \subset I_1 \subset I_2 \subset . . . $ - последовательность идеалов в K, то $\exists \ m \ \in N : \forall \ n \geq m :I_n = I_m$.
    \item [2] Любой идеал в K является конечнопорожденным.
    \item[3] Если $(a_0) \subset (a_0, a_1) \subset (a_0, a_1, a_2) ... $ - последовательность идеалов в K, то $\exists m \in \mathbb{N} \ \forall \ n \geq m: \ a_n \in (a_1, ..., a_m)$
\end{itemize}
Рассмотрим наши кольца:
\begin{itemize}
    \item [] \textbf{КГИ.} В КГИ любой идеалел конечнопорожден, а значит, выполняется второе определение
     \item [] \textbf{$Z[\sqrt{d}]$ (d - любое целое число).}
\end{itemize}